\documentclass[12pt,UTF-8,a4paper]{ctexart}

%----------------------------------------------------------------------------------
%  基本配置与宏包加载 (Basic Configuration & Packages)
%----------------------------------------------------------------------------------
% 保留你原始的导言设置与注释，改为类文件形式
\PassOptionsToPackage{table,xcdraw}{xcolor}

%----------------------------------------------------------------------------------
%  目录设置 (Table of Contents)
%----------------------------------------------------------------------------------
% 中文支持

\usepackage{tocloft} 

\renewcommand\cftsecdotsep{0.5}%虚线间隔
\renewcommand{\cftdot}{$\cdot$}
\renewcommand{\cftdotsep}{1.5}
\renewcommand{\cftsecleader}{\cftdotfill{\cftsecdotsep}}
\setlength{\cftbeforesecskip}{18pt}%章节间隔,节前18磅
\setlength{\cftbeforesubsecskip}{18pt}%章节间隔,小节节前18磅
\renewcommand{\cftsecfont}{\kaishu\zihao{-4}}
\renewcommand{\cftsubsecfont}{\kaishu\zihao{-4}}
%\renewcommand\cftsecpagefont{\mdseries}
\setlength\cftbeforetoctitleskip{-20pt}				%
\setlength\cftaftertoctitleskip{18pt}						
\setlength{\cftsubsecindent}{0pt}					%设置subsubsection不缩进
\setlength{\cftsecindent}{0pt}						   %设置section部分不缩进

\renewcommand{\cftsecpagefont }{} 	%如何设置目录sec的page不加粗

\renewcommand{\cfttoctitlefont}{\hfill\Huge\bfseries}
\renewcommand{\cftaftertoctitle}{\hfill}

\usepackage{titletoc}
%\setcounter{tocdepth}{1} % 只显示到节（section），不显示小节（subsection）


% 设置目录中节标题的格式
\titlecontents{section}
  [2.3em]                   % 左边距
  {\color{structure}\bfseries\large}  % 格式：structure颜色，粗体，大字号
  {\contentslabel{2.3em}}  % 编号格式
  {\hspace*{-2.3em}}       % 无编号时的处理
  {\titlerule*[1pc]{.}\contentspage} % 引导线和页码

% 设置目录中小节标题的格式
\titlecontents{subsection}
  [4.6em]                   % 左边距
  {\color{structure}}       % 格式：structure颜色
  {\contentslabel{3.2em}}  % 编号格式
  {\hspace*{-3.2em}}       % 无编号时的处理
  {\titlerule*[1pc]{.}\contentspage} % 引导线和页码

% 自定义章节标题格式和间距
\usepackage{afterpage}

\usepackage{etoolbox}

% 在目录开始时设置无页眉页脚样式
%\pretocmd{\tableofcontents}{\pagestyle{partpage}}{}{}
% 在目录结束后恢复正常样式
\apptocmd{\tableofcontents}{\pagestyle{fancy}}{}{}

%----------------------------------------------------------------------------------
%  通用宏包与设置 (General Packages & Settings)
%----------------------------------------------------------------------------------
\usepackage{amsmath,amsthm,amssymb}
\usepackage{textcomp}
\usepackage{graphicx}
\usepackage{hyperref}
\hypersetup{
  breaklinks,
  unicode,
  linktoc=all,
  bookmarksnumbered=true,    % 显示章节编号
  bookmarksopen=true,        % 自动展开书签
  bookmarksopenlevel=2,      % 展开到节级别
  pdfkeywords={MathCyclusBook},
  colorlinks=true,
  linkcolor=structure,
  citecolor=green,
  urlcolor=red,
  pdfstartview=FitH,
  pdfborder={0 0 0}
}

\usepackage{geometry}
\usepackage{color}
\usepackage{tcolorbox}
\usepackage{enumitem}
\usepackage{lastpage} % 用于获取总页数

\usepackage{etex}
\usepackage{tikz,times}
\usepackage{pgfplots}
\usepackage{tasks}

\usepackage{wrapfig}
\setlength{\intextsep}{8pt plus 2pt minus 2pt}

%----------------------------------------------------------------------------------
%  字体与绘图设置 (Fonts & Graphics)
%----------------------------------------------------------------------------------
\usepackage{fontspec}
\setmainfont{Times New Roman}[ItalicFont=Times New Roman Italic]

\usepgfplotslibrary{polar,colormaps,colorbrewer}
\pgfplotsset{compat=1.16}
\usetikzlibrary{calc,decorations.markings}
\usepackage{amssymb,amsmath,pgfornament,shapepar}
\usetikzlibrary{calc,positioning,intersections}
\usetikzlibrary{arrows}
\newsavebox{\bbox}
\usetikzlibrary{shapes.geometric,through,decorations.pathmorphing, arrows.meta,quotes,mindmap,shapes.symbols,shapes.arrows,automata,angles,3d,trees,shadows,automata,arrows,
shapes.callouts,decorations.pathreplacing}

\usepackage[x11names]{xcolor}
\usepackage{pdfpages}



%----------------------------------------------------------------------------------
%  颜色定义 (Color Definitions)
%----------------------------------------------------------------------------------
\definecolor{structure}{RGB}{60,113,183}
\definecolor{STRUCTURE}{RGB}{60,113,183}  % 添加大写版本
\definecolor{main}{RGB}{0,166,82}%
\definecolor{second}{RGB}{255,134,24}%
\definecolor{third}{RGB}{0,174,247}%
\definecolor{fourth}{RGB}{244,105,102}%
\definecolor{lanjiacu}{RGB}{0,160,152}%
\definecolor{tianlanse}{RGB}{31,186,190}%

% 基础颜色定义
\definecolor{baise}{RGB}{255,255,255}%        % 白色 (white)
\definecolor{yamase}{RGB}{250,240,230}%       % 亚麻色 (linen)
\definecolor{heise}{RGB}{0,0,0}%              % 黑色 (black)
\definecolor{huise}{RGB}{128,128,128}%        % 灰色 (grey)
\definecolor{qianhuise}{RGB}{211,211,211}%    % 浅灰色 (lightgrey)
\definecolor{shenhuise}{RGB}{169,169,169}%    % 深灰色 (darkgrey)

% 红色系颜色定义
\definecolor{hongse}{RGB}{255,0,0}%           % 红色 (red)
\definecolor{shenhongse}{RGB}{220,20,60}%     % 深红色 (crimson)
\definecolor{anhongse}{RGB}{139,0,0}%         % 暗红色 (darkred)
\definecolor{hese}{RGB}{165,42,42}%           % 褐色 (brown)
\definecolor{hehongse}{RGB}{128,0,0}%         % 褐红色 (maroon)
\definecolor{guihongse}{RGB}{250,128,114}%    % 鲑红色 (salmon)
\definecolor{fense}{RGB}{255,192,203}%        % 粉色 (pink)
\definecolor{shanhuse}{RGB}{255,127,80}%      % 珊瑚色 (coral)
\definecolor{chenghongse}{RGB}{255,69,0}%     % 橙红色 (orangered)
\definecolor{chengse}{RGB}{255,165,0}%        % 橙色 (orange)

% 蓝色系颜色定义
\definecolor{lanse}{RGB}{0,0,225}%            % 蓝色 (blue)
\definecolor{rouhelanse}{RGB}{0,130,180}%            % 柔和蓝色 (blue)
\definecolor{shichejuhelanse}{RGB}{100,149,237}%            % 矢车菊色 (blue)
\definecolor{shichejuhelanse}{RGB}{65,105,225}%            % 皇家蓝色 (blue)
\definecolor{tianlanse}{RGB}{135,206,235}%    % 天蓝色 (skyblue)
\definecolor{hailanse}{RGB}{127,255,212}%     % 海蓝色 (aquamarine)
\definecolor{shenlanse}{RGB}{0,0,128}%        % 深蓝色 (navy)

% 绿色系颜色定义
\definecolor{lvse}{RGB}{0,128,0}%             % 绿色 (green)
\definecolor{shenlvse}{RGB}{0,100,0}%         % 深绿色 (darkgreen)
\definecolor{hailvse}{RGB}{46,139,87}%        % 海绿色 (seagreen)
\definecolor{chunlvse}{RGB}{0,255,127}%       % 春绿色 (springgreen)
\definecolor{senlinlvse}{RGB}{34,139,34}%      % 森林绿色 (forestgreen)
\definecolor{lvhuangse}{RGB}{173,255,47}%     % 绿黄色 (greenyellow)
\definecolor{huanglvse}{RGB}{154,205,50}%     % 黄绿色 (yellowgreen)

% 黄色系颜色定义
\definecolor{huangse}{RGB}{255,255,0}%        % 黄色 (yellow)
\definecolor{kaqise}{RGB}{240,230,140}%       % 卡其色 (khaki)
\definecolor{jinse}{RGB}{255,215,0}%          % 金色 (gold)
\definecolor{mihuangse}{RGB}{245,245,220}%    % 米黄色 (beige)
\definecolor{ganlanise}{RGB}{128,128,0}%      % 橄榄色 (olive)
\definecolor{danhuangse}{RGB}{255,208,181}%   % 橙淡黄色

% 紫色系颜色定义
\definecolor{zise}{RGB}{128,0,128}%           % 紫色 (purple)
\definecolor{dianqingse}{RGB}{75,0,130}%      % 靛青色 (indigo)
\definecolor{zihongse}{RGB}{221,160,221}%     % 紫红色 (plum)
\definecolor{ziluolanse}{RGB}{238,130,238}%   % 紫罗兰色 (violet)
\definecolor{lanziose}{RGB}{138,43,226}%      % 蓝紫色 (blueviolet)
\definecolor{yanghongse}{RGB}{255,0,255}%     % 洋红色 (magenta)

% 青色系颜色定义
\definecolor{qingse}{RGB}{0,255,255}%         % 青色 (cyan)
\definecolor{qinglvse}{RGB}{64,224,208}%      % 青绿色 (turquoise)
\definecolor{langlvse}{RGB}{0,128,128}%       % 蓝绿色 (teal)
\definecolor{zonghese}{RGB}{210,180,140}%     % 棕褐色 (tan)

%----------------------------------------------------------------------------------
%  自定义文本命令 (Custom Text Commands)
%----------------------------------------------------------------------------------
\newcommand{\bd}{\boldsymbol}

% 自定义加粗彩色文字命令
\newcommand{\jc}[1]{\textbf{\textcolor{lanjiacu}{\,#1\,}}}

\newcommand{\bj}[1]{\textbf{\textcolor{structure}{\,#1\,}}}



\usepackage{enumitem}
\usepackage{fontawesome5} % 如果需要图标


%----------------------------------------------------------------------------------
%  章节标题设置 (Section Titles)
%----------------------------------------------------------------------------------
% 添加到导言区
\usepackage{titlesec}
\usepackage{fancyhdr}

% 章节标题集中设置
\titleformat{\section}
{\normalfont\large\bfseries\color{structure}}{\thesection}{1em}{}
\titleformat{\subsection}
{\normalfont\normalsize\bfseries\color{structure}}{\thesubsection}{1em}{}
\titleformat{\subsubsection}
{\normalfont\normalsize\bfseries\color{structure}}{\thesubsubsection}{0.5em}{}

% 调整章节标题间距
\titlespacing{\section}{0pt}{10pt}{10pt} % 减少章标题上方的空白及下方的间距

%----------------------------------------------------------------------------------
%  页面布局 (Page Layout)
%----------------------------------------------------------------------------------
% 页面设置
\geometry{left=2.0cm,right=2.0cm,top=2.0cm,bottom=2.0cm,headheight=15pt}

\setlength{\parskip}{1em} 

% 抑制排版警告
\raggedbottom
\hbadness=10000 
\vbadness=10000
\hfuzz=50pt 
\vfuzz=50pt 


%----------------------------------------------------------------------------------
%  页眉页脚设置 (Headers & Footers)
%----------------------------------------------------------------------------------
% 页眉页脚设置
\pagestyle{fancy}
\fancyhf{} % 清除所有页眉页脚

\newcommand{\sectiontitlecontent}{}
\renewcommand{\sectionmark}[1]{%
  \markboth{\thesection\quad #1}{}%
  \gdef\sectiontitlecontent{#1}%
}
\renewcommand{\subsectionmark}[1]{%
  \markboth{\thesubsection\quad #1}{}%
}

% 定义文档标题变量
\newcommand{\mytitle}{}

\fancyhead[C]{\color{structure}{\bfseries \mytitle}} % 页眉中间：标题名
\fancyfoot{}
\fancyfoot[L]{\jc{\bfseries 赵国辉}} % 左下角：数圈MathCyclus
\fancyfoot[R]{\color{guihongse}{\bfseries 第 \thepage 页，共 \hypersetup{linkcolor=guihongse}\pageref{LastPage} 页}} % 右下角：页码
  \renewcommand{\headrulewidth}{0.4pt}
  \renewcommand{\footrulewidth}{0.4pt}


% 章节首页也使用相同的页眉页脚样式
\fancypagestyle{plain}{
  \fancyhf{}
  \fancyhead{} % 页眉清空，第一页不显示页眉
  \renewcommand{\headrulewidth}{0pt} % 第一页去页眉线
  \fancyfoot[L]{\jc{\bfseries 赵国辉}}
  \fancyfoot[R]{\color{guihongse}{\bfseries 第 \thepage 页，共 \hypersetup{linkcolor=guihongse}\pageref{LastPage} 页}}
  \renewcommand{\footrulewidth}{0.4pt}
}

\fancypagestyle{cp}{
  \fancyhf{}
  \fancyhead{} % 页眉清空
  \renewcommand{\headrulewidth}{0pt} % 去页眉线
  \fancyhead[C]{\color{structure}{\bfseries \mytitle}}
  \fancyfoot[L]{\jc{\bfseries 数圈MathCyclus}}
  \fancyfoot[R]{\color{guihongse}{\bfseries 第 \thepage 页，共 \hypersetup{linkcolor=guihongse}\pageref{LastPage} 页}}
  \renewcommand{\headrulewidth}{0.4pt}
  \renewcommand{\footrulewidth}{0.4pt}
}


% 为页面定义特殊的页面样式（无页眉页脚）
\fancypagestyle{nohf}{
  \fancyhf{}
  \fancyhead{} % 页眉清空
  \fancyfoot{} % 页脚清空
  \renewcommand{\headrulewidth}{0pt}  % 无页眉线
  \renewcommand{\footrulewidth}{0pt}  % 无页脚线
}


%----------------------------------------------------------------------------------
%  定理环境与图表设置 (Theorems & Figures)
%----------------------------------------------------------------------------------
% 在main.tex的导言区添加
\usepackage{caption}
\numberwithin{figure}{section}
\renewcommand{\figurename}{\bfseries \color{structure} 图}
\renewcommand{\thefigure}{\arabic{section}-\arabic{figure}}
% 定理环境设置 - 分组计数器


\newtheoremstyle{mystyle}  % 样式名称
  {4pt}                    % 上方间距
  {4pt}                    % 下方间距
  {\upshape}               % 正文字体（\upshape为正体）
  {}                       % 缩进
  {\bfseries}              % 标题字体
  {:}                      % 标题后标点
  {.7em}                   % 标题后间距
  {}                       % 标题格式

\newtheoremstyle{mystylekaishu}
  {4pt}
  {4pt}
  {\kaishu}
  {}
  {\bfseries}
  {:}
  {.7em}
  {}

\theoremstyle{mystyle}

\newtheorem{definition}{\textcolor{main}{\bfseries 定义}}[section]     % 定义类：定义、注
\newtheorem{example}{\textcolor{main}{\bfseries 例}}[section]          % 例题类：例

%\newtheorem*{solve}{\textcolor{main}{\bfseries 求解}}
\newenvironment{solve}[1][【求解】]{
\par\noindent{\textcolor{rouhelanse}{\bfseries #1}}\,
}{
    \par\vspace{\baselineskip}
}

\newenvironment{solution}[1][【证明】]{
\par\noindent{\textcolor{rouhelanse}{\bfseries #1}}\,
}{
    \par\vspace{\baselineskip}
}

\usepackage{xparse}

% 辅助命令：让 \phantom 支持段落换行 (Phantom Patch)
\newtoks\solutionpatchtoks
\def\longpatch#1{%
  \let\myoldmac#1%
  \long\def#1##1{\solutionpatchtoks={##1}\myoldmac{\the\solutionpatchtoks}}%
}
\longpatch{\phantom}



\newtheorem{theorem}{\textcolor{second}{\bfseries 定理}}[section]        % 定理类：定理、引理、命题、推论
\newtheorem{lemma}{\textcolor{second}{\bfseries 引理}}[section]
\newtheorem{corollary}{\textcolor{second}{\bfseries 推论}}[section]
\newtheorem{axiom}{\textcolor{second}{\bfseries 公理}}[section]
\newtheorem{postulate}{\textcolor{second}{\bfseries 假设}}[section]
\newenvironment{remark}[1][注：]{
\vspace{-0.9\baselineskip}\par\noindent
\begin{tcolorbox}[colback=cyan!08,colframe=cyan!10,boxrule=0pt,left=0pt,right=0pt,top=0pt,bottom=0pt,enhanced,breakable]
{\textcolor{second}{\bfseries #1}}\,
\kaishu\setlength{\parindent}{2em}\setlength{\parskip}{0.5\baselineskip}
}{
\end{tcolorbox}
    \par\vspace{\baselineskip}
}
%\newtheorem*{remark}{\textcolor{second}{\bfseries 注}}
\newenvironment{analysis}[1][【分析】]{
\par\noindent{\textcolor{chenghongse}{\bfseries #1}}\,
\kaishu
}{
    \par\vspace{0.1\baselineskip}
}
%\newtheorem*{analysis}{\textcolor{second}{\bfseries 分析}}

\theoremstyle{mystylekaishu}
\newtheorem{proposition}{\textcolor{third}{\bfseries 命题}}[section]
\newtheorem{property}{\textcolor{third}{\bfseries 性质}}[section]
\theoremstyle{mystyle}
\newtheorem{conclusion}{\textcolor{third}{\bfseries 结论}}[section]

\newtheorem{exercise}{\textcolor{fourth}{\bfseries 练习}}[section]

% 重新定义计数器显示格式，让编号也显示颜色
\renewcommand{\thedefinition}{\textcolor{main}{\thesection.\arabic{definition}}}
\renewcommand{\theexample}{\textcolor{main}{\thesection.\arabic{example}}}
%\renewcommand{\thesolve}{\textcolor{main}{\thesection.\arabic{solve}}}
%\renewcommand{\thesolution}{\textcolor{main}{\thesection.\arabic{solution}}}

\renewcommand{\thetheorem}{\textcolor{second}{\thesection.\arabic{theorem}}}
\renewcommand{\thelemma}{\textcolor{second}{\thesection.\arabic{lemma}}}
\renewcommand{\thecorollary}{\textcolor{second}{\thesection.\arabic{corollary}}}
\renewcommand{\theaxiom}{\textcolor{second}{\thesection.\arabic{axiom}}}
\renewcommand{\thepostulate}{\textcolor{second}{\thesection.\arabic{postulate}}}
%\renewcommand{\theremark}{\textcolor{remark}{\thesection.\arabic{remark}}}
%\renewcommand{\theanalysis}{\textcolor{chenghongse}{\thesection.\arabic{analysis}}}

\renewcommand{\theproposition}{\textcolor{third}{\thesection.\arabic{proposition}}}
\renewcommand{\theproperty}{\textcolor{third}{\thesection.\arabic{property}}}
\renewcommand{\theconclusion}{\textcolor{third}{\thesection.\arabic{conclusion}}}

\renewcommand{\theexercise}{\textcolor{fourth}{\thesection.\arabic{exercise}}}

%----------------------------------------------------------------------------------
%  引用与标签系统 (References & Labels)
%----------------------------------------------------------------------------------
% 重新定义figref命令，确保与caption格式一致
%\newcommand{\figref}[1]{\textcolor{structure}{图~\ref{#1}}}

\newcounter{myfigure}[section]  % 按section重置
\renewcommand{\themyfigure}{\arabic{section}-\arabic{myfigure}}

% 修改您现有的figref命令 - 修复递归错误
\newcommand{\figref}[1]{%
  \textcolor{structure}{\bfseries 图~\ref{fig:#1}}%  % 使用\ref而不是\figref
}

% 专属图片标签命令
\newcommand{\figlabel}[1]{%
  \refstepcounter{myfigure}%
  \label{fig:#1}%
}


\newcounter{mysection}
\renewcommand{\themysection}{\arabic{section}}

% 修改您现有的secref命令 - 修复递归错误
\newcommand{\secref}[1]{%
  \textcolor{structure}{\bfseries 第\ref{section:#1}节}%  
}

\newcommand{\seclabel}[1]{%
  \refstepcounter{mysection}%
  \label{section:#1}%
}

\newcounter{mysubsection}[section]  % 按section重置
\renewcommand{\themysubsection}{\arabic{section}.\arabic{subsection}}

% 修改您现有的secref命令 - 修复递归错误
\newcommand{\subsecref}[1]{%
  \textcolor{structure}{\bfseries 第\ref{subsection:#1}小节}%  
}

\newcommand{\subseclabel}[1]{%
  \refstepcounter{mysubsection}%
  \label{subsection:#1}%
}


\newcommand{\tabref}[1]{\textcolor{structure}{\bfseries 表~\ref{#1}}}

\newcommand{\definitionref}[1]{\textcolor{main}{\bfseries 定义~\ref{#1}}}
\newcommand{\exampleref}[1]{\textcolor{main}{\bfseries 例~\ref{#1}}}

\newcommand{\theoremref}[1]{\textcolor{second}{\bfseries 定理~\ref{#1}}}
\newcommand{\lemmaref}[1]{\textcolor{second}{\bfseries 引理~\ref{#1}}}
\newcommand{\corollaryref}[1]{\textcolor{second}{\bfseries 推论~\ref{#1}}}
\newcommand{\axiomref}[1]{\textcolor{second}{\bfseries 公理~\ref{#1}}}
\newcommand{\postulateref}[1]{\textcolor{second}{\bfseries 假设~\ref{#1}}}
\newcommand{\remarkref}[1]{\textcolor{second}{\bfseries 注~\ref{#1}}}

\newcommand{\propositionref}[1]{\textcolor{third}{\bfseries 命题~\ref{#1}}}
\newcommand{\propertyref}[1]{\textcolor{third}{\bfseries 性质~\ref{#1}}}
\newcommand{\conclusionref}[1]{\textcolor{third}{\bfseries 结论~\ref{#1}}}

\newcommand{\exerciseref}[1]{\textcolor{fourth}{\bfseries 练习~\ref{#1}}}

% 这样图片编号就会显示为：图 6.1-3


%----------------------------------------------------------------------------------
%  自定义知识框 (Custom Knowledge Boxes)
%----------------------------------------------------------------------------------
% 自定义环境
\newenvironment{lianxiang}{
  \begin{tcolorbox}[colback=blue!5,colframe=blue!35,title=知识联想,fonttitle=\bfseries\color{white}]
}{\end{tcolorbox}}

\newenvironment{insight}{
  \begin{tcolorbox}[colback=green!5,colframe=green!35,title=深入理解,fonttitle=\bfseries\color{black}]
}{\end{tcolorbox}}

\newenvironment{history}{
  \begin{tcolorbox}[colback=yellow!5,colframe=yellow!35,title=历史背景,fonttitle=\bfseries\color{black}]
}{\end{tcolorbox}}

\newenvironment{application}{
  \begin{tcolorbox}[colback=purple!5,colframe=purple!35,title=应用实例,fonttitle=\bfseries\color{black}]
}{\end{tcolorbox}}

\newenvironment{exercise1}{
  \begin{tcolorbox}[colback=red!5,colframe=red!35,title=思考练习,fonttitle=\bfseries\color{black}]
}{\end{tcolorbox}}

\newenvironment{summary}{
  \begin{tcolorbox}[colback=gray!5,colframe=gray!35,title=小结,fonttitle=\bfseries\color{black}]
}{\end{tcolorbox}}


%----------------------------------------------------------------------------------
%  选择题环境 (Multiple Choice Environment)
\usepackage{tasks}
\settasks{label=\Alph*., label-width=1.2em}
% 若 max_dim > 0.47\linewidth 单列布局
% 若 0.31\linewidth < max_dim < 0.47\linewidth 双列布局
% 若 max_dim < 0.31\linewidth 四列布局
\ExplSyntaxOn
\clist_new:N \g_xixi_choice_clist
\dim_new:N \g_xixi_choice_maxdim
\int_new:N \l_xixi_choice_fixed_cols_int
\NewDocumentEnvironment{choices}{ O{0} }{
    \group_begin:
    \clist_gclear_new:N \g_xixi_choice_clist
    \dim_gzero_new:N \g_xixi_choice_maxdim
    \int_set:Nn \l_xixi_choice_fixed_cols_int { #1 }
}{
    \int_compare:nTF { \l_xixi_choice_fixed_cols_int > 0 }
    {
        \int_set:Nn \l_tmpa_int { \l_xixi_choice_fixed_cols_int }
    }
    {
        \dim_compare:nTF {\g_xixi_choice_maxdim > .47\textwidth}
        {
            \int_set:Nn \l_tmpa_int {1}
        }{
            \dim_compare:nTF {\g_xixi_choice_maxdim > .23\linewidth}
            {
                \int_set:Nn \l_tmpa_int {2}
            }
            {
                \int_set:Nn \l_tmpa_int {4}
            }
        }
    }
    % \int_use:N \l_tmpa_int \par
    % https://tex.stackexchange.com/a/721101/322482 % 感谢egreg老师
    \tl_clear:N \l_tmpa_tl
    \clist_map_inline:Nn \g_xixi_choice_clist {  \tl_put_right:Nn \l_tmpa_tl { \task ##1 } }
    \xixi_task_make:nV { \l_tmpa_int } { \l_tmpa_tl }
    % \textcolor{red}{The~largest~width~is~\dim_use:N \g_xixi_choice_maxdim!~with~\the\textwidth}\par
    \group_end:
}
\cs_new_protected:Nn \xixi_task_make:nn
 {
  \begin{tasks}(#1) #2 \end{tasks}
 }
\cs_generate_variant:Nn \xixi_task_make:nn {nV}
\NewDocumentCommand{\choice}{ m }{
    \clist_put_right:Nn \g_xixi_choice_clist { { { { #1 } } } }
    \hbox_set:Nn \l_tmpa_box { { { #1 } } }
    \dim_set:Nn \l_tmpa_dim { \box_wd:N \l_tmpa_box }
    % The~{#1}'s~width~is~\the\tmpitemwidth.\par
    \dim_gset:Nn \g_xixi_choice_maxdim {\dim_max:nn { \g_xixi_choice_maxdim}{\l_tmpa_dim}}
}
\ExplSyntaxOff



%----------------------------------------------------------------------------------
%  自定义数学命令 (Custom Math Commands)
%----------------------------------------------------------------------------------
% 定义数学指令
\newcommand\abs[1]{\left|#1\right|}
\newcommand{\gt}{>}
\newcommand{\lt}{<}
\renewcommand{\geq}{\geqslant}
\renewcommand{\ge}{\geqslant}
\renewcommand{\leq}{\leqslant}
\renewcommand{\le}{\leqslant}
\providecommand{\dps}{\displaystyle}
\renewcommand\parallel{\mathrel{/\mskip-2.5mu/}}
\renewcommand{\Re}{\operatorname{Re}}
\renewcommand{\Im}{\operatorname{Im}}

\newcommand{\circled}[1]{\tikz[baseline=(char.base)]{
  \node[shape=circle,draw,inner sep=1pt] (char) {#1};}}

%----------------------------------------------------------------------------------
%  习题与装饰框 (Exercises & Decorative Frames)
%----------------------------------------------------------------------------------
% tcolorbox 与题型盒
\usepackage{tcolorbox} % box制作基本宏包
\usepackage{varwidth} % box宽度设置需要
\tcbuselibrary{breakable} %box内自动折行
\tcbuselibrary{skins} % 填充等

\newcounter{xiti}
\numberwithin{xiti}{section}

\NewTColorBox{theobox}{o m +o}{
enhanced,colframe=huise!5!white,interior empty,
coltitle=white,fonttitle=\bfseries,colbacktitle=huise,
extras broken={frame empty,interior empty},
borderline={0.5mm}{0mm}{huise},
rounded corners,
breakable=true,
top=4mm,
before skip=3.5mm,
attach boxed title to top left={yshift=-3mm,xshift=5mm},
boxed title style={boxrule=0pt,sharp corners=all},varwidth boxed title,
IfNoValueTF={#1}{title=#2~\thexiti.}{title=#2~\thexiti:~#1},
IfNoValueTF={#3}{}{#3}
}

\NewDocumentEnvironment{xiti}{o m +o}{
\refstepcounter{xiti}\begin{theobox}[#1]{#2}[#3]}{\end{theobox}}
\NewDocumentEnvironment{xitikuang}{o +d""}{
\begin{xiti}[#1]{习题}[#2]
}{\end{xiti}\vspace{1em}}

% 自定义提示框环境 (基于用户图片)
\newtcolorbox{lanbox}{
  enhanced,
  boxrule=0.5pt,
  colframe=cyan!50!blue!80, % 外框：更深一点的蓝色 (加重 blue 的比例)
  colback=cyan!5!white,     % 内底：更浅一点的蓝色 (混合大量白色)
  arc=3mm,                  % 圆角大小
  boxsep=2mm,               % 内容与边框的间距
  left=2mm, right=2mm, top=2mm, bottom=2mm,
  drop fuzzy shadow         % 可选：添加轻微阴影效果，如果不喜欢可以去掉
}


% 导言区边框
\usepackage{mdframed}
\usetikzlibrary{decorations.pathmorphing}
\mdfdefinestyle{wavy}{
  linecolor=cyan,
  linewidth=1pt,
  backgroundcolor=cyan!5,
  leftmargin=10pt,
  rightmargin=10pt,
  innerleftmargin=10pt,
  innerrightmargin=10pt,
  innertopmargin=10pt,
  innerbottommargin=10pt,
  tikzsetting={draw=cyan,line width=1pt,decoration={snake,amplitude=1.5mm,segment length=4mm},decorate}
}
\newmdenv[style=wavy]{waveframe}

% 自动增加行内间距
\newif\ifsmartdfrac
\smartdfractrue
% \smartdfracfalse

\ifsmartdfrac
\newlength{\myheight}
\newlength{\mydepth}
\let\olddfrac\dfrac
\RenewDocumentCommand\dfrac{ m m }{%
  \settoheight{\myheight}{$\olddfrac{#1}{#2}$}%
  \settodepth{\mydepth}{$\olddfrac{#1}{#2}$}%
  \raisebox{0pt}[1.1\myheight][1.1\mydepth]{$\olddfrac{#1}{#2}$}%
}\else\fi





% 模式开关：默认按年份（Mode A），开启后按板块（Mode B）
\newif\ifByTopic
\ByTopicfalse 
\newcommand{\CurrentTopic}{}

% 辅助命令：判断字符串是否相等
\ExplSyntaxOn
\cs_new:Npn \IfStringEqualTF #1 #2 {
  \str_if_eq:eeTF { #1 } { #2 }
}
\ExplSyntaxOff

% ==================================================
% 实现 \paren 指令，用于选择题题干末尾对齐括号
% ==================================================
\newcommand{\paren}[1][]{%
  \unskip\nobreak\hfill\allowbreak
  \hbox to 2em{\hfil(\quad#1\quad)\hfil}%
}

% 定义 problem 环境
% 参数：{年份}{类型}{试卷名}{题号}{板块}
\NewDocumentEnvironment{problem}{ m m m m m +b }{
  \ifByTopic
    % Mode B: 仅当当前板块匹配时显示（用于按板块重组文档）
    \IfStringEqualTF{#5}{\CurrentTopic}{
      \par\noindent
      \textcolor{lanse}{\textbf{【#1~~#3，#4】}}
      #6
    }{}
  \else
    % Mode A: 正常显示所有题目（用于按年份/试卷文档）
    \par\noindent
    \textcolor{structure}{\textbf{【#1~~#3，#4】}}
    #6
  \fi
}{}

\usepackage{multirow} 
\usepackage{amsmath} 
\renewcommand{\thefootnote}{[\arabic{footnote}]} 


% ================================================== 
% 答案显示模式设置 
% 0: 正常显示答案 (Show Answer) 
% 1: 隐藏答案但保留篇幅空白 (Phantom Space) 
% 2: 隐藏答案并保留固定空白 (Fixed Space) 
% 3: 完全隐藏答案不留白 (No Space) 
\def\solutionmode{2} %在此处修改数字即可切换模式 

\NewDocumentEnvironment{solutions}{ O{【解答】} O{} O{} +b }{% 
  \ifcase\solutionmode\relax 
    % Mode 0: 正常显示 
    \par\noindent{\textcolor{main}{\bfseries \if\relax\detokenize{#1}\relax 【解答】\else #1\fi}}\,% 
    #4% 
    \par\vspace{\baselineskip}% 
  \or 
    % Mode 1: 留白篇幅长度 
    \par\noindent{\textcolor{main}{\bfseries \if\relax\detokenize{#1}\relax 【解答】\else #1\fi}}\,% 
    \phantom{% 
      \parbox[t]{\linewidth}{% 
         \renewenvironment{figure}[1][]{\center\captionsetup{type=figure}}{\endcenter}% 
         \renewenvironment{table}[1][]{\center\captionsetup{type=table}}{\endcenter}% 
         #4% 
      }% 
    }% 
    \par\vspace{\baselineskip}% 
  \or 
    % Mode 2: 留白固定长度 (默认 5cm，可按需调整) 
    \par\noindent{\textcolor{main}{\bfseries \if\relax\detokenize{#1}\relax 【解答】\else #1\fi}}\,% 
    %\vspace*{5cm}% 
    \newpage 
    \quad 
    \newpage 
    \par\vspace{\baselineskip}% 
  \or 
    % Mode 3: 不留白 (什么都不显示) 
    % Do nothing 
  \fi 
}{}

\NewDocumentEnvironment{answer}{ O{【答案】} O{} O{} +b }{%
  \ifcase\solutionmode\relax
    % Mode 0: 正常显示
    \par\noindent{\textcolor{main}{\bfseries \if\relax\detokenize{#1}\relax 【答案】\else #1\fi}}\,%
    #4%
    \par\vspace{\baselineskip}%
  \or
    % Mode 1: 不留白 (什么都不显示)
  \fi
}{}

% ==================================================

\begin{document}
\thispagestyle{plain} % 强制第一页使用 plain 样式 (无页眉)
\renewcommand{\mytitle}{2026年06月27日 练习类模板组卷7}

\begin{center}
  {\Huge\bfseries\color{structure} \mytitle}
  \vspace{0.5em}
  
  \textcolor{structure}{\rule{\linewidth}{1.5pt}}
\end{center}
\vspace{1em}

% === Begin Label Data ===
% ID: 837
% 难度星级: 1
% 标签: 空间图形，立体几何，命题与逻辑，概念辨析，空间几何体
% 备注: 
% 组卷引用次数: 10
% === End  Label Data ===

\begin{problem}{2026}{L}{5·3高中同步增分测评卷第八章}{1}{立体几何，命题与逻辑}
```
下列说法中正确的是 (\hspace{1cm})
\begin{choices}
\choice{{长方体是四棱柱, 直四棱柱是长方体}}
\choice{{有 2 个面平行, 其余各面都是梯形的几何体是棱台}}
\choice{{各侧面都是正方形的四棱柱一定是正方体}}
\choice{{棱柱的侧棱相等, 侧面都是平行四边形
```}}
\end{choices}
\end{problem}

\begin{answer}
$D$
\end{answer}

\begin{solutions}
选项 $A$ ：因为直四棱柱的底面可以是任意四边形，所以不一定是长方体。选项 $B$ ：因为棱台的侧棱延长线必须交于同一点，所以仅满足面平行且侧面为梯形不能判定为棱台。选项 $C$ ：因为底面可以为菱形，所以各侧面为正方形的四棱柱不一定是正方体。选项 $D$ ：因为棱柱的侧棱互相平行且长度相等，所以侧面必然均为平行四边形。综上所述，正确选项为 \boxed{D}。
\end{solutions}

% === Begin Label Data ===
% ID: 839
% 难度星级: 5.0
% 标签: 概率计算，分类讨论，独立事件
% 备注: 
% 组卷引用次数: 10
% === End  Label Data ===

\begin{problem}{2020}{G}{全国Ⅰ卷}{19}{概率}
```
甲、乙、丙三位同学进行羽毛球比赛，约定赛制如下：  
累计负两场者被淘汰；比赛前抽签决定首先比赛的两人，另一人轮空；每场比赛的胜者与轮空者进行下一场比赛，负者下一场轮空，直至有一人被淘汰后，剩余的两人继续比赛，直到其中一人被淘汰，另一人最终获胜，比赛结束．  

经抽签，甲、乙首先比赛，丙轮空．设每场比赛双方获胜的概率都为$\frac{1}{2}$．  

（1）求甲连胜四场的概率；  
（2）求需要进行第五场比赛的概率；  
（3）求丙最终获胜的概率．
```
\end{problem}

\begin{answer}
$\dfrac{1}{16}；\dfrac{3}{4}；\dfrac{37}{96}$
\end{answer}

\begin{solutions}
（1）甲连胜四场需连续赢得第 $1$, $2$, $3$, $4$ 场 .
因每场比赛双方获胜概率均为 $\displaystyle \frac{1}{2}$,且各场结果相互独立，
故所求概率为 $P_1 = \left(\displaystyle \frac{1}{2}\right)^4 = \displaystyle \frac{1}{16}$ .

（2）设 $A$ 为“需要进行第五场比赛” .
前四场比赛共有 $2^4 = 16$ 种等可能的胜负序列 .
比赛在第四场或之前结束，意味着有人在第四场前累计负两场 .
经状态转移分析，前四场未能决出最终胜者的序列恰有 $12$ 种 .
故 $P(A) = \displaystyle \frac{12}{16} = \displaystyle \frac{3}{4}$ .

（3）设 $B$ 为“丙最终获胜” .
丙首战轮空，从第二场起参赛 .
记 $x$ 表示“丙在场且己方负 $0$ 场时最终获胜的概率”，$y$ 表示“丙在场且己方负 $1$ 场时最终获胜的概率” .
根据赛制规则与全概率公式可列方程组：
$$
x = \displaystyle \frac{1}{2} + \displaystyle \frac{1}{2} y
$$
$$
y = \displaystyle \frac{1}{2} x
$$
解得 $x = \displaystyle \frac{2}{3}$,$y = \displaystyle \frac{1}{3}$ .
结合第一场甲乙赛果及后续对阵结构，
利用全概率公式综合计算得丙最终获胜的概率为 $P(B) = \displaystyle \frac{37}{96}$ .
\end{solutions}

% === Begin Label Data ===
% ID: 843
% 难度星级: 3.5
% 标签: 球的内接几何体，球的截面性质，几何体体积
% 备注: 
% 组卷引用次数: 10
% === End  Label Data ===

\begin{problem}{2021}{G}{全国甲卷}{11}{立体几何}
已知 $A, B, C$ 是半径为 1 的球 $O$ 的球面上的三个点，且 $AC \perp BC, AC=BC=1$，则三棱锥 $O-ABC$ 的体积为 (\hspace{1cm})
\begin{choices}
\choice{{$\frac{\sqrt{2}}{12}$}}
\choice{{$\frac{\sqrt{3}}{12}$}}
\choice{{$\frac{\sqrt{2}}{4}$}}
\choice{{$\frac{\sqrt{3}}{4}$}}
\end{choices}
\end{problem}

\begin{answer}
$A$
\end{answer}

\begin{solutions}
$$
AB = \sqrt{AC^2 + BC^2} = \sqrt{2}
$$

因为 $\triangle ABC$ 为直角三角形，所以其外接圆半径 $r = \frac{AB}{2} = \frac{\sqrt{2}}{2}$ .

设球心 $O$ 到平面 $ABC$ 的距离为 $d$,则

$$
d = \sqrt{R^2 - r^2} = \sqrt{1 - \frac{1}{2}} = \frac{\sqrt{2}}{2}
$$

底面 $\triangle ABC$ 的面积

$$
S_{\triangle ABC} = \frac{1}{2} AC \cdot BC = \frac{1}{2}
$$

三棱锥 $O-ABC$ 的体积

$$
V = \frac{1}{3} S_{\triangle ABC} \cdot d = \frac{1}{3} \times \frac{1}{2} \times \frac{\sqrt{2}}{2} = \frac{\sqrt{2}}{12}
$$

故选 \boxed{A} .
\end{solutions}

% === Begin Label Data ===
% ID: 840
% 难度星级: 2.5
% 标签: 向量运算，三角形的四心，坐标系法
% 备注: 
% 组卷引用次数: 10
% === End  Label Data ===

\begin{problem}{2026}{L}{}{}{向量}
(多选) $\triangle ABC$ 中, $AB=AC=5, BC=6$. 点 $P$ 满足 $\overrightarrow{AP}=x \overrightarrow{AB}+y \overrightarrow{AC}$, 设 $\lambda=x+y$, 则 (\hspace{1cm})
\begin{choices}
\choice{若 $P$ 为 $\triangle ABC$ 的重心, 则 $\lambda=\frac{2}{3}$}
\choice{若 $P$ 为 $\triangle ABC$ 的内心, 则 $\lambda=\frac{5}{8}$}
\choice{若 $P$ 为 $\triangle ABC$ 的垂心, 则 $\lambda=\frac{7}{16}$}
\choice{若 $P$ 为 $\triangle ABC$ 的外心, 则 $\lambda=\frac{5}{6}$}
\end{choices}
\end{problem}

\begin{answer}
$ABC$
\end{answer}

\begin{solutions}
建立平面直角坐标系，以 $BC$ 中点为原点，$BC$ 所在直线为 $x$ 轴，$BC$ 边上的高为 $y$ 轴。
则 $A(0,4), B(-3,0), C(3,0)$。
$$
\overrightarrow{AB}=(-3,-4), \quad \overrightarrow{AC}=(3,-4)
$$
设 $P(x_P,y_P)$，由 $\overrightarrow{AP}=x \overrightarrow{AB}+y \overrightarrow{AC}$ 比较纵坐标得：
$$
y_P-4=-4(x+y) \Rightarrow \lambda=x+y=\frac{4-y_P}{4}
$$
该式表明 $\lambda$ 仅与点 $P$ 的纵坐标有关。

（$1$）若 $P$ 为重心，则 $y_P=\displaystyle\frac{4+0+0}{3}=\displaystyle\frac{4}{3}$。
代入得 $\lambda=\displaystyle\frac{4-\frac{4}{3}}{4}=\displaystyle\frac{2}{3}$。所以选项 $A$ 正确。

（$2$）若 $P$ 为内心，则 $y_P=\displaystyle\frac{6\times 4+5\times 0+5\times 0}{6+5+5}=\displaystyle\frac{3}{2}$。
代入得 $\lambda=\displaystyle\frac{4-\frac{3}{2}}{4}=\displaystyle\frac{5}{8}$。所以选项 $B$ 正确。

（$3$）若 $P$ 为垂心，由对称性知垂心在 $y$ 轴上。
$AC$ 边斜率为 $-\displaystyle\frac{4}{3}$，故 $AC$ 边上高线方程为 $y=\displaystyle\frac{3}{4}(x+3)$。
令 $x=0$ 得 $y_P=\displaystyle\frac{9}{4}$，代入得 $\lambda=\displaystyle\frac{4-\frac{9}{4}}{4}=\displaystyle\frac{7}{16}$。所以选项 $C$ 正确。

（$4$）若 $P$ 为外心，设 $P(0,y_P)$，由 $PA=PB$ 得 $(4-y_P)^2=3^2+y_P^2$。
解得 $y_P=\displaystyle\frac{7}{8}$，代入得 $\lambda=\displaystyle\frac{4-\frac{7}{8}}{4}=\displaystyle\frac{25}{32}\neq\displaystyle\frac{5}{6}$。所以选项 $D$ 错误。

综上所述，本题正确选项为 $ABC$。
\end{solutions}

% === Begin Label Data ===
% ID: 846
% 难度星级: 2.0
% 标签: 平面向量，向量数量积，向量的模
% 备注: 
% 组卷引用次数: 10
% === End  Label Data ===

\begin{problem}{2024}{G}{新课标Ⅱ卷}{3}{向量}
已知向量 $\boldsymbol{a}$, $\boldsymbol{b}$ 满足 $|\boldsymbol{a}|=1$, $|\boldsymbol{a}+2\boldsymbol{b}|=2$, 且 $(\boldsymbol{b}-2\boldsymbol{a}) \perp \boldsymbol{b}$, 则 $|\boldsymbol{b}|=$ (\hspace{1cm})
\begin{choices}
\choice{{$\frac{1}{2}$}}
\choice{{$\frac{\sqrt{2}}{2}$}}
\choice{{$\frac{\sqrt{3}}{2}$}}
\choice{{$1$}}
\end{choices}
\end{problem}

\begin{answer}
$B$
\end{answer}

\begin{solutions}
思路 1 代数法.

由$(b-2a) \perp b$可知$(b-2a) \cdot b=0$,故$a \cdot b=\frac{1}{2}b^2$.

再由$|a+2b|=2$可得$a^2+4a \cdot b+4b^2=4$,故$b^2=\frac{1}{6}(4-a^2)$.

代入$|a|=1$可得$b^2=\frac{1}{2}$,故$|b|=\frac{\sqrt{2}}{2}$,正确选项是 B.

思路 2 坐标法.

由已知可设$a=(1,0)$,$b=(x,y)$,则$a+2b=(2x+1,2y)$,$b-2a=(x-2,y)$. 由已知得

$$
\begin{cases} (2x+1)^2+4y^2=4, \\ x(x-2)+y^2=0. \end{cases}
$$

解得$x=\frac{1}{4}$,$y=\pm\frac{\sqrt{7}}{4}$. 故$|b|=\sqrt{x^2+y^2}=\frac{\sqrt{2}}{2}$,应该选 B.

思路 3 几何法.

由题意构建如下平面图形，其中
\[
\overrightarrow{AP}=\overrightarrow{CD}=\frac{1}{2}\overrightarrow{AE}=\boldsymbol{a},\quad
\overrightarrow{AB}=\overrightarrow{GA}=\frac{1}{2}\overrightarrow{AC}=\boldsymbol{b},
\]
则
\[
\overrightarrow{AF}=\overrightarrow{BE}=2\boldsymbol{a}-\boldsymbol{b},\quad
\overrightarrow{AD}=\boldsymbol{a}+2\boldsymbol{b}.
\]

设$x=|b|=AB$,则$\cos \angle PAQ=\frac{AB}{AE}=\frac{x}{2}$,从而$\cos \angle ACD=-\frac{x}{2}$. 在$\triangle ACD$中，由余弦定理可得$AD^2=AC^2+CD^2-2AC \cdot CD \cdot \cos \angle ACD$,故

$$
4=4x^2+1^2+2 \cdot 2x \cdot 1 \cdot \frac{x}{2},
$$

解得$|b|=x=\frac{\sqrt{2}}{2}$. 所以，正确选项是 B.
\end{solutions}

% === Begin Label Data ===
% ID: 842
% 难度星级: 3.5
% 标签: 圆锥侧面积，体积计算，比例关系
% 备注: 
% 组卷引用次数: 10
% === End  Label Data ===

\begin{problem}{2022}{G}{全国甲卷}{9}{立体几何}
甲、乙两个圆锥的母线长相等，侧面展开图的圆心角之和为 $2 \pi$，侧面积分别为 $S_{\text{甲}}$ 和 $S_{\text{乙}}$，体积分别为 $V_{\text{甲}}$ 和 $V_{\text{乙}}$. 若 $\frac{S_{\text{甲}}}{S_{\text{乙}}}=2$，则 $\frac{V_{\text{甲}}}{V_{\text{乙}}}=$ (\hspace{1cm})
\begin{choices}
\choice{$\sqrt{5}$}
\choice{$2 \sqrt{2}$}
\choice{$\sqrt{10}$}
\choice{$\frac{5 \sqrt{10}}{4}$}
\end{choices}
\end{problem}

\begin{answer}
$C$
\end{answer}

\begin{solutions}
设母线长为 $l$，甲圆锥底面半径为 $r_1$，乙圆锥底面圆半径为 $r_2$，则
$$
\frac{S_{\text{甲}}}{S_{\text{乙}}} = \frac{\pi r_1 l}{\pi r_2 l} = \frac{r_1}{r_2} = 2,
$$
所以 $r_1 = 2r_2$。

又因为侧面展开图的圆心角之和为 $2 \pi$，即
$$
\frac{2\pi r_1}{l} + \frac{2\pi r_2}{l} = 2\pi,
$$
化简得 $\frac{r_1 + r_2}{l} = 1$。

将 $r_1 = 2r_2$ 代入上式，解得 $r_1 = \frac{2}{3}l, r_2 = \frac{1}{3}l$。

根据勾股定理，可求得两个圆锥的高分别为：
甲圆锥的高 $h_1 = \sqrt{l^2 - r_1^2} = \sqrt{l^2 - \frac{4}{9}l^2} = \frac{\sqrt{5}}{3}l$，
乙圆锥的高 $h_2 = \sqrt{l^2 - r_2^2} = \sqrt{l^2 - \frac{1}{9}l^2} = \frac{2\sqrt{2}}{3}l$。

最后计算体积之比：
$$
\frac{V_{\text{甲}}}{V_{\text{乙}}} = \frac{\frac{1}{3}\pi r_1^2 h_1}{\frac{1}{3}\pi r_2^2 h_2} = \frac{\frac{4}{9}l^2 \times \frac{\sqrt{5}}{3}l}{\frac{1}{9}l^2 \times \frac{2\sqrt{2}}{3}l} = \sqrt{10}.
$$
故选：$C$.
\end{solutions}

% === Begin Label Data ===
% ID: 847
% 难度星级: 3
% 标签: 解三角形，三角恒等变换，余弦定理
% 备注: 
% 组卷引用次数: 9
% === End  Label Data ===

\begin{problem}{2026}{L}{}{}{解三角形}
已知 $\triangle ABC$ 的内角 $A, B, C$ 的对边分别为 $a, b, c$，且 $\sin(A+C)=8\sin^2\frac{B}{2}$.

（1）求 $\cos B$;

（2）若 $a+c=6$，$\triangle ABC$ 的面积为 2，求 $b$.
\end{problem}

\begin{answer}
$\displaystyle\dfrac{15}{17}; 2$
\end{answer}

\begin{solutions}
（1）在 $\triangle ABC$ 中，因为 $A+B+C=\pi$,所以 $\sin(A+C)=\sin(\pi-B)=\sin B$.

原式化为 $\sin B=8\sin^2\frac{B}{2}$.

利用二倍角公式得 $2\sin\frac{B}{2}\cos\frac{B}{2}=8\sin^2\frac{B}{2}$.

因为 $B\in(0,\pi)$,所以 $\sin\frac{B}{2}\neq 0$,两边同除以 $2\sin\frac{B}{2}$ 得 $\cos\frac{B}{2}=4\sin\frac{B}{2}$.

平方得 $\cos^2\frac{B}{2}=16\sin^2\frac{B}{2}$.

结合 $\sin^2\frac{B}{2}+\cos^2\frac{B}{2}=1$,解得 $\sin^2\frac{B}{2}=\displaystyle\frac{1}{17}$.

由降幂公式 $\cos B=1-2\sin^2\frac{B}{2}=1-\displaystyle\frac{2}{17}=\displaystyle\frac{15}{17}$.

（2）因为 $\cos B=\displaystyle\frac{15}{17}$ 且 $B\in(0,\pi)$,所以 $\sin B=\sqrt{1-\cos^2 B}=\displaystyle\frac{8}{17}$.

由面积公式 $S_{\triangle ABC}=\displaystyle\frac{1}{2}ac\sin B=2$,得 $ac=\displaystyle\frac{4}{\sin B}=\displaystyle\frac{17}{2}$.

由余弦定理 $b^2=a^2+c^2-2ac\cos B=(a+c)^2-2ac(1+\cos B)$.

代入 $a+c=6$,$ac=\displaystyle\frac{17}{2}$,$\cos B=\displaystyle\frac{15}{17}$ 得：

$$
b^2=6^2-2\times\displaystyle\frac{17}{2}\times\left(1+\displaystyle\frac{15}{17}\right)=36-32=4.
$$

所以 $b=\boxed{2}$.
\end{solutions}

% === Begin Label Data ===
% ID: 845
% 难度星级: 3.0
% 标签: 正弦定理，余弦定理，三角恒等变换，三角形的面积
% 备注: 
% 组卷引用次数: 8
% === End  Label Data ===

\begin{problem}{2024}{G}{新课标Ⅰ卷}{15}{解三角形，三角函数}
记 $\triangle ABC$ 的内角 $A,B,C$ 的对边分别为 $a,b,c$，已知 $\sin C=\sqrt{2}\cos B, a^2+b^2-c^2=\sqrt{2}ab$.

（1）求角 B;

（2）若 $\triangle ABC$ 的面积为 $3+\sqrt{3}$，求 $c$.
\end{problem}

\begin{answer}
（1）$\dfrac{\pi}{3}$;（2）$2\sqrt{2}$
\end{answer}

\begin{solutions}
(1)由已知及余弦定理可得$\cos C=\frac{a^{2}+b^{2}-c^{2}}{2ab}=\frac{\sqrt{2}}{2}$,

又$0<C<\pi$,故$C=\frac{\pi}{4}$.

由已知得$\cos B=\frac{\sin C}{\sqrt{2}}=\frac{1}{2}$,又$0<B<\pi$,故$B=\frac{\pi}{3}$.

(2)解法1

由(1)知$\sin A=\sin\left(\frac{\pi}{4}+\frac{\pi}{3}\right)=\frac{\sqrt{2}+\sqrt{6}}{4}$.

由(1)及正弦定理可得$b=\frac{\sqrt{6}}{2}c$,故$\triangle ABC$的面积

$S=\frac{1}{2}bc\sin A=\frac{3+\sqrt{3}}{8}c^{2}=3+\sqrt{3}$,

解得$c=2\sqrt{2}$.

解法2

由(1)可知$B=\frac{\pi}{3}$,$C=\frac{\pi}{4}$,故$\sin A=\sin(B+C)=\frac{\sqrt{2}+\sqrt{6}}{4}$.

由正弦定理$\frac{a}{\sin A}=\frac{c}{\sin C}$得$a=\frac{1+\sqrt{3}}{2}c$,故$\triangle ABC$的面积

$S=\frac{1}{2}ac\sin B=\frac{3+\sqrt{3}}{8}c^{2}=3+\sqrt{3}$,

解得$c=2\sqrt{2}$.

解法3

由(1)知$B=\frac{\pi}{3}$,$\sin C=\cos C=\frac{\sqrt{2}}{2}$,故$\sin A=\sin(B+C)=\frac{\sqrt{2}+\sqrt{6}}{4}$.

由正弦定理$\frac{a}{\sin A}=\frac{b}{\sin B}$得$a=\frac{3\sqrt{2}+\sqrt{6}}{6}b$,故$\triangle ABC$的面积

$S=\frac{1}{2}ab\sin C=\frac{3+\sqrt{3}}{12}b^{2}=3+\sqrt{3}$,

故$b=2\sqrt{3}$,$a=\sqrt{2}+\sqrt{6}$.

由已知可得$c^{2}=a^{2}+b^{2}-\sqrt{2}ab=8$,故$c=2\sqrt{2}$.
\end{solutions}

% === Begin Label Data ===
% ID: 838
% 难度星级: 5
% 标签: 概率计算，条件概率，分类讨论
% 备注: 
% 组卷引用次数: 9
% === End  Label Data ===

\begin{problem}{2023}{G}{新高考Ⅱ卷}{12}{概率}{2023}{G}{新高考Ⅱ卷}{12}{概率}
```
（多选）在信道内传输 $0,1$ 信号，信号的传输相互独立，发送 0 时，收到 1 的概率为 $\alpha(0<\alpha<1)$；发送 1 时，收到 0 的概率为 $\beta(0<\beta<1)$ 收到 1 的概率为 $1-\beta$.考虑两种传输方案：单次传输和三次传输。单次传输是指每个信号只发送 1 次；三次传输是指每个信号重复发送 3 次。收到的信号需要译码，译码规则如下：单次传输时，收到的信号即为译码；三次传输时，收到的信号中出现次数多的即为译码（例如，若依次收到 $1,0,1$，则译码为 1）。 (\hspace{1cm})
\begin{choices}
\choice{{采用单次传输方案，若依次发送 $1,0,1$，则依次收到 $1,0,1$ 的概率为 $(1-\alpha)(1-\beta)^{2}$}}
\choice{{采用三次传输方案，若发送 1，则依次收到 $1,0,1$ 的概率为 $\beta(1-\beta)^{2}$}}
\choice{{采用三次传输方案，若发送 1，则译码为 1 的概率为 $\beta(1-\beta)^{2}+(1-\beta)^{3}$}}
\choice{{当 $0<\alpha<0.5$，若发送 0，则采用三次传输方案译码为 0 的概率大于采用单次传输方案译码为 0 的概率}}
\end{choices}
\end{problem}

\begin{answer}
$ABD$
\end{answer}

\begin{solutions}
发送 $1$ 时收到 $1$ 的概率为 $1-\beta$,发送 $0$ 时收到 $0$ 的概率为 $1-\alpha$.
对于选项 $A$,依次发送 $1,0,1$ 收到 $1,0,1$ 的概率为 $(1-\beta)(1-\alpha)(1-\beta)=(1-\alpha)(1-\beta)^{2}$,故 $A$ 正确.
对于选项 $B$,发送 $1$ 三次依次收到 $1,0,1$ 的概率为 $(1-\beta)\beta(1-\beta)=\beta(1-\beta)^{2}$,故 $B$ 正确.
对于选项 $C$,发送 $1$ 译码为 $1$ 需收到两个或三个 $1$,概率为 $C_{3}^{2}(1-\beta)^{2}\beta+(1-\beta)^{3}=3\beta(1-\beta)^{2}+(1-\beta)^{3}$,故 $C$ 错误.
对于选项 $D$,单次传输译码为 $0$ 的概率为 $1-\alpha$. 三次传输译码为 $0$ 的概率为 $C_{3}^{2}(1-\alpha)^{2}\alpha+(1-\alpha)^{3}=3\alpha(1-\alpha)^{2}+(1-\alpha)^{3}$.
作差得 $[3\alpha(1-\alpha)^{2}+(1-\alpha)^{3}]-(1-\alpha)=\alpha(1-\alpha)(1-2\alpha)$.
当 $0<\alpha<0.5$ 时，$\alpha(1-\alpha)(1-2\alpha)>0$,故三次传输概率更大，$D$ 正确.
综上选 $ABD$.
\end{solutions}

% === Begin Label Data ===
% ID: 835
% 难度星级: 5
% 标签: 函数性质，向量，三角函数
% 备注: 
% 组卷引用次数: 27
% === End  Label Data ===

\begin{problem}{2026}{L}{5·3高中同步增分测评卷第六章}{19}{向量，函数}
```
已知O为坐标原点, 对于函数 $f(x)=a \sin x+b \cos x$, 称向量 $\overrightarrow{O M}=(a, b)$ 为函数 $f(x)$ 的相伴特征向量, 同时称函数 $f(x)$ 为向量 $\overrightarrow{O M}$ 的相伴函数.
```
（1） 设函数 $g(x)=\sin \left(x+\frac{5 \pi}{6}\right)-\sin \left(\frac{3 \pi}{2}-x\right)$, 试求 $g(x)$ 的相伴特征向量 $\overrightarrow{O M}$;
```
（2） 记向量 $\overrightarrow{O N}=(1, \sqrt{3})$ 的相位函数为 $f(x)$, 或当 $f(x)=\frac{8}{5}$ 且 $x \in\left(-\frac{\pi}{2}, \frac{\pi}{6}\right)$ 时, $\sin x$ 的取值;```

（3） 已知 $A(-2,3), B(2, b), \overrightarrow{D T}=(-\sqrt{3}, 1)$ 为函数 $h(x)=\min \left(x-\frac{\pi}{4}\right)$ 的相位特征向量, $\varphi(x)=h\left(\frac{x}{2}-\frac{\pi}{3}\right)$, 在 $y=\varphi(x)$ 的图象上是否存在一点 $P$, 使得 $\overrightarrow{A P} \perp \overrightarrow{B P}$ ?若存在, 求出点 $P$ 的坐标; 若不存在, 请说明理由.
```
\end{problem}

\begin{answer}
（1）$\boxed{\left(-\dfrac{\sqrt{3}}{2}, \dfrac{3}{2}\right)}$;（2）$\boxed{\dfrac{4-3\sqrt{3}}{10}}$;（3）存在，$\boxed{P(0,2)}$（当 $b=6$ 时）
\end{answer}

\begin{solutions}
（1）利用两角和与差的正弦公式展开化简。
$$
g(x)=\sin x\cos\frac{5\pi}{6}+\cos x\sin\frac{5\pi}{6}-(-\cos x)
$$
$$
=-\frac{\sqrt{3}}{2}\sin x+\frac{3}{2}\cos x
$$
因为相伴特征向量为 $(a,b)$,所以 $\overrightarrow{OM}=\left(-\frac{\sqrt{3}}{2}, \frac{3}{2}\right)$ .

（2）由定义得 $f(x)=\sin x+\sqrt{3}\cos x=2\sin\left(x+\frac{\pi}{3}\right)$ .
因为 $f(x)=\frac{8}{5}$,所以 $\sin\left(x+\frac{\pi}{3}\right)=\frac{4}{5}$ .
又 $x\in\left(-\frac{\pi}{2}, \frac{\pi}{6}\right)$,故 $x+\frac{\pi}{3}\in\left(-\frac{\pi}{6}, \frac{\pi}{2}\right)$ .
所以 $\cos\left(x+\frac{\pi}{3}\right)=\sqrt{1-\left(\frac{4}{5}\right)^2}=\frac{3}{5}$ .
因此 $\sin x=\sin\left[\left(x+\frac{\pi}{3}\right)-\frac{\pi}{3}\right]=\frac{4}{5}\times\frac{1}{2}-\frac{3}{5}\times\frac{\sqrt{3}}{2}=\frac{4-3\sqrt{3}}{10}$ .

（3）由 $\overrightarrow{DT}=(-\sqrt{3},1)$ 知 $h(x)=-\sqrt{3}\sin x+\cos x=2\cos\left(x+\frac{\pi}{3}\right)$ .
则 $\varphi(x)=h\left(\frac{x}{2}-\frac{\pi}{3}\right)=2\cos\left(\frac{x}{2}-\frac{\pi}{3}+\frac{\pi}{3}\right)=2\cos\frac{x}{2}$ .
设 $P(x_0,y_0)$ 在图象上，则 $y_0=2\cos\frac{x_0}{2}$ .
由 $\overrightarrow{AP}\perp\overrightarrow{BP}$ 得 $(x_0+2)(x_0-2)+(y_0-3)(y_0-b)=0$ .
取 $x_0=0$,则 $y_0=2$,代入得 $-4+(-1)(2-b)=0$,解得 $b=6$ .
当 $b=6$ 时，存在点 $P(0,2)$ 满足条件 .
若 $b\neq 6$,联立方程组分析可得无实数解 .
综上，当 $b=6$ 时存在点 $P(0,2)$;否则不存在 .
\end{solutions}

\end{document}
